\scanpage{276}
- \mspage{127} -
\margin{IX}
10.5. \uline{Réduction de} 10.4 \uline{à un cas particulier}.

10.5.1. L'unicité d'un \(u\) satisfaisant la condition énoncée (même pour un seul \(\ell\)) est claire. D'autre part, l'existence signifie simplement que chaque \(u_\ell\) applique en fait \(\uline{M}\otimes\uline{M}'\) dans \(\unsure{\mathbb Z_S}\) et que l'homomorphisme ainsi obtenu \(\uline{M}\otimes\uline{M}'\longrightarrow\mathbb Z_S\) est indépendant de \(\ell.\) En fait, il suffit même de prouver l'assertion analogue pour les \(u_\ell\otimes_{\mathbb Z_\ell}\mathbb Q_\ell\) et un \(u_{\mathbb Q}:\uline{M}_{\mathbb Q}\otimes\uline{M}'_{\mathbb Q}\longrightarrow\mathbb Q_S.\)

Supposons alors qu'on ait établi l'existence après changement de base \(S'\longrightarrow S\) comme dans 9.3. Alors il résulte immédiatement de (10.3.5) qu'il existe un homomorphisme
\[
u:\uline{M}\otimes\uline{M}'\longrightarrow\frac{1}{e(S'/S)}\mathbb Z_S\subset\mathbb Q_S
\]
compatible avec les \(u_\ell\otimes_{\mathbb Z_\ell}\mathbb Q_\ell,\) d'où a). Quant à b), il est trivial alors que s'il est vrai pour \(A_{K'},\) il l'est aussi pour \(A_K.\)

Par suite, nous pouvons, pour prouver 10.4, nous permettre n'importe quel changement de traits \(S'\longrightarrow S\) dominant. Ceci nous permet en particulier de supposer \(S\) complet à corps résiduel algébriquement clos. Dans ce cas, \(\uline{M}\) et \(\uline{M}'\) sont des schémas en groupes constants, de valeur \(M,M'\) des modules libres de type fini sur \(\mathbb Z,\) et a) et b) se reformulent avantageusement en termes de ces derniers.

10.5.2. Supposons maintenant que \(A_K\) soit isomorphe à un facteur direct d'un schéma abélien \(\bar A_K\) sur \(k\) à réduction semi-stable, et que 10.4 soit prouvé pour ce dernier. Alors il en résulte aussitôt que 10.4 est valable également pour \(A_K,\) compte tenu du comportement fonctoriel évident des réseaux \(\uline{M},\uline{M}',\) et des accouplements de monodromie, vis-à-vis des sommes
\figure{\begin{tikzpicture}\draw (0,0)--(.25,0);\draw[->] (1,-.4)--(.9,-.55);\draw[->] (3,-.4)--(2.85,-.5);\draw[thick] (2,-1)--(2.5,-1);\end{tikzpicture}}
\margin{445}
\margin{276}
\margin{445}
